% © 2026 Ing. David Jaroš — CC BY-NC-ND 4.0
%
% This work is licensed under a Creative Commons Attribution-NonCommercial-NoDerivatives
% 4.0 International License (CC BY-NC-ND 4.0).
%
% License History: Earlier drafts (up to v0.3) were released under CC BY 4.0.
% From v0.4 onward, all material is released under CC BY-NC-ND 4.0 to protect
% the integrity of the theoretical work during ongoing academic development.

\documentclass[11pt,a4paper]{article}
\usepackage{amsmath, amssymb, amsthm}
\usepackage{hyperref}
\usepackage{geometry}
\geometry{margin=2.5cm}

\newtheorem{theorem}{Theorem}
\newtheorem{proposition}{Proposition}
\newtheorem{definition}{Definition}
\newtheorem{remark}{Remark}
\newtheorem{conjecture}{Conjecture}

\title{Mirror Sector Dynamics at \texorpdfstring{$p = 139$}{p=139}\\
\large Twin-Prime Structure, Set~B Hecke Eigenvalues, and Physical Interpretation}
\author{Ing.\ David Jaroš}
\date{2026}

\begin{document}
\maketitle

\begin{abstract}
We document and analyse the ``mirror sector'' numerical observation in Unified
Biquaternion Theory (UBT): a second set of modular newforms (Set~B) whose
Hecke eigenvalues at the twin prime $p = 139$ reproduce the charged lepton
mass ratios to similar accuracy as Set~A at $p = 137$.  We examine the
theoretical mechanisms that could explain the mirror-symmetry between the
$p = 137$ and $p = 139$ sectors, including parity symmetry of the KK tower,
CPT duality, and modular symmetry of the UBT partition function.  The term
``mirror sector'' is used descriptively for a mathematically observed pattern;
no physical mirror world is claimed.
All claims are labelled with their proof status.
\end{abstract}

\tableofcontents

% ======================================================================
\section{The Set~B Numerical Observation}
\label{sec:setb}
% ======================================================================

\begin{remark}[Epistemic notice]
  All mass-ratio agreements in this document are numerical observations.
  No theoretical derivation connecting Hecke eigenvalues to physical masses
  is known.  See the full epistemic statement in
  \texttt{research/mirror\_sector\_modular\_status.md}.
\end{remark}

\subsection{Set~B: Hecke eigenvalues at $p = 139$}

From \texttt{research/mirror\_sector\_modular\_status.md} §2:

\begin{center}
\begin{tabular}{lccccc}
\hline
Form & Level $N$ & Weight $k$ & $a_{139}$ & Ratio & Target & Error \\
\hline
$g_1$ (electron ref.) & $N_1'$ & 2 & $a_{139}^{(1)}$ & 1 & 1 & — \\
$g_2$ (muon) & $N_2'$ & 4 & $a_{139}^{(2)}$ & $\approx 206.77$ & 206.768 & $<0.1\%$ \\
$g_3$ (tau) & $N_3'$ & 6 & $a_{139}^{(3)}$ & $\approx 3477$ & 3477.23 & $<0.2\%$ \\
\hline
\end{tabular}
\end{center}

\textbf{Status: [NO]} — numerical observation.  Source: \texttt{reports/hecke\_lepton/mirror\_world\_139.md}.

% ======================================================================
\section{Theoretical Mechanisms for Mirror Symmetry}
\label{sec:mechanisms}
% ======================================================================

\subsection{Parity of the KK tower}

The KK modes $n$ and $-n$ are related by $\psi \to -\psi$ (parity of the
imaginary time direction).  In UBT, $\psi$-parity is a symmetry of the
classical action (Status: \textbf{[L0]}).

\begin{conjecture}[Mirror sector from $\psi$-parity]\label{conj:psi_parity}
  The twin-prime pair $(p, p+2) = (137, 139)$ corresponds to modes $+n$ and
  $-n$ of the $\psi$-circle, related by $\psi$-parity.  The mirror sector at
  $p = 139$ is the $\psi$-parity image of the physical sector at $p = 137$.

  \textbf{Status: [L2]} — motivated by the parity structure, but the
  connection to twin primes is not explained.
\end{conjecture}

\subsection{CPT duality}

In a CPT-symmetric theory, the spectrum of particles and antiparticles are
related.  If the $p = 139$ sector corresponds to antiparticles, then the
mirror observation reflects CPT duality.

\begin{remark}[Status]
  CPT symmetry in UBT follows from the Hermitian structure of the
  $\mathbb{H}\otimes\mathbb{C}$ algebra (Status: \textbf{[L1]}).
  However, CPT relates particles to antiparticles \emph{of the same mass},
  not to a sector with a different prime $p$.  The CPT interpretation of
  the mirror sector therefore requires a specific mechanism connecting
  prime shifts to CPT operations.  Status: \textbf{[O]}.
\end{remark}

\subsection{Modular Atkin--Lehner involution}

Every newform $f$ of level $N$ has an Atkin--Lehner involution:
\begin{equation}
\label{eq:al}
  W_N f = \varepsilon_f f, \quad \varepsilon_f = \pm 1.
\end{equation}

Under $W_N$, the prime $p$ is mapped to $p' = N/p \pmod{N}$ (for $p | N$).

\begin{conjecture}[Mirror sector from Atkin--Lehner]\label{conj:al}
  The Set~B forms at $p = 139$ are the Atkin--Lehner images of the Set~A
  forms at $p = 137$ under $W_N$.

  \textbf{Status: [L2]}.  This would require the levels $N_1', N_2', N_3'$
  to satisfy $N'/137 \equiv 139 \pmod{N'}$, which imposes arithmetic
  constraints that have not been checked.
\end{conjecture}

% ======================================================================
\section{Physical Interpretation Candidates}
\label{sec:physical}
% ======================================================================

\subsection{No physical mirror world claimed}

The term ``mirror sector'' is used solely as a descriptive label for the
Set~B numerical observation.  No physical mirror-world scenario (mirror
SM, mirror dark matter, etc.) is claimed.

\subsection{Possible structural roles}

\begin{enumerate}
  \item \textbf{Functional equation}: The $L$-functions $L(f_k, s)$ and
    $L(g_k, s)$ may be related by a functional equation connecting
    $p = 137$ and $p = 139$ evaluations.  This would be a purely arithmetic
    statement with no direct physical content.

  \item \textbf{Vacuum degeneracy}: If the UBT vacuum is doubly degenerate
    (two inequivalent vacua at $p = 137$ and $p = 139$), then Set~A and
    Set~B would correspond to two different vacuum sectors.  The physical
    universe occupies one of them.  Status: \textbf{[L2]}.

  \item \textbf{Numerical coincidence}: The match at both $p = 137$ and
    $p = 139$ may be a numerical coincidence within the expected accuracy
    of the modular-form approach.  Status: must be tested by computing
    the probability of two independent primes both matching.
\end{enumerate}

% ======================================================================
\section{Vacuum Stability Analysis}
\label{sec:vacuum}
% ======================================================================

Following \texttt{consolidation\_project/mirror\_sector/vacuum\_stability.tex}:

\subsection{Effective potential at $p = 137$ and $p = 139$}

If the effective potential $V_{\mathrm{eff}}(n)$ is evaluated at the
prime $p$ corresponding to each sector, the difference
$\Delta V = V_{\mathrm{eff}}(137) - V_{\mathrm{eff}}(139)$ determines
which vacuum is preferred.

\begin{remark}[Status]
  The computation of $V_{\mathrm{eff}}$ at prime values has not been
  performed from the UBT action.  The mirror-sector analysis in
  \texttt{consolidation\_project/mirror\_sector/vacuum\_stability.tex}
  uses a phenomenological ansatz for $V_{\mathrm{eff}}$.
  Status: \textbf{[L2]}.
\end{remark}

% ======================================================================
\section{Summary and Open Problems}
\label{sec:open}
% ======================================================================

\begin{center}
\begin{tabular}{lll}
\hline
Claim & Status & Source \\
\hline
Set~B matches at $p=139$ & \textbf{[NO]} numerical & \texttt{reports/hecke\_lepton/mirror\_world\_139.md} \\
$\psi$-parity mechanism & \textbf{[L2]} conjecture & This document \\
Atkin--Lehner mechanism & \textbf{[L2]} conjecture & This document \\
Physical mirror world & \textbf{NOT CLAIMED} & This document \\
$\Delta V$ between sectors & \textbf{[O]} open & \texttt{mirror\_sector/vacuum\_stability.tex} \\
\hline
\end{tabular}
\end{center}

\textbf{Key open problem}: No theoretical mechanism connecting the twin-prime
pair $(137, 139)$ to the $\psi \to -\psi$ symmetry or any other UBT symmetry
has been derived.  The observation remains numerically interesting but
theoretically unexplained.

% ======================================================================
\section{Cross-References}
\label{sec:xref}
% ======================================================================

\begin{tabular}{ll}
\hline
Topic & File \\
\hline
Mirror sector overview & \texttt{research/mirror\_sector\_modular\_status.md} \\
Set~A Hecke analysis & \texttt{research/hecke\_generation\_structure.tex} \\
SageMath computation & \texttt{research\_tracks/three\_generations/step5\_hecke\_search\_results.tex} \\
Vacuum stability & \texttt{consolidation\_project/mirror\_sector/vacuum\_stability.tex} \\
Statistical significance & \texttt{reports/hecke\_lepton/statistical\_significance.md} \\
\hline
\end{tabular}

\end{document}
